\loai{Xác định tên hạt nhân}
\begin{vidu}%[3P7-4.2B][Loai1]
Trong phản ứng sau đây: 
\begin{align*}
n+\Hn{92}{235}{U}\to \Hn{42}{95}{Mo}+\Hn{57}{139}{La}+2\text{X}+7\beta^-
\end{align*} 
Hạt $X$ là
\choice
{electron}
{\True neutron}
{proton}
{heli}
\loigiai{
Áp dụng định luật bảo toàn số khối và bảo toàn điện tích:
\begin{align*}
\left\{\begin{aligned}& 1+235=95+139+2A+7.0  \\
& 0+92=42+57+2Z-7.
\end{aligned}\right.
\Rightarrow \left\{\begin{aligned}& A=1 \\
& Z=0 
\end{aligned}\right.\Rightarrow\Hn{0}{1}{n}.
\end{align*}
}
\end{vidu}
\begin{vidu}%[3P7-4.2B][Loai1]
\nguon{CĐ - 2013} Trong phản ứng hạt nhân: $\Hn{9}{19}{F}+p\rightarrow \Hn{8}{16}{O}+X$, hạt X là
\choice
{electron}
{positron}
{proton}
{\True hạt $\alpha$}
\loigiai{\begin{itemize}
\item ĐLBT số khối: $19+1=16+A\Rightarrow A = 4$
\item ĐLBT điện tích: $9+1=8+Z \Rightarrow Z = 2$. 
\item Vậy X là hạt $\alpha$.
\end{itemize}
}
\end{vidu}

\loai{Xác định số proton, nuclôn, neutron}
\begin{vidu}%[3P7-4.2Y][Loai2]
Cho phản ứng hạt nhân: $\Hn{17}{35}{Cl}+ \Hn{Z}{A}{X} \to n+\Hn{18}{37}{Ar}$. Trong đó hạt X có
\choice
{\True $Z=1$; $A=3$}
{$Z=2$; $A=4$}
{$Z=2$; $A=3$}
{$Z=1$; $A=1$}
\loigiai{
Áp dụng định luật bảo toàn số khối và bảo toàn điện tích:
\begin{align*}
\left\{\begin{aligned}& 35+A=1+37 \\
& 17+Z=0+18 
\end{aligned}\right.
\Rightarrow \left\{\begin{aligned}& A=3 \\
& Z=1 
\end{aligned}\right.
\end{align*}
}
\end{vidu}

\begin{vidu}%[3P7-4.2B][Loai2]
\nguon{QG - 2017} Cho phản ứng hạt nhân: $\Hn{2}{4}{He}+\Hn{7}{14}{N} \to \Hn{1}{1}{H}+X$. Số proton và neutron của hạt nhân $X$ lần lượt là 	
\choice
{\True 8 và 9}
{9 và 17}
{9 và 8}
{8 và 17}
\loigiai{ 
\begin{align*}
\text{Ta có:}\left\{\begin{aligned}& 4+14=1+A \\
& 2+7=1+Z 
\end{aligned}\right.\Rightarrow \left\{\begin{aligned}& A=17 \\
& Z=8 
\end{aligned}\right.\Rightarrow N=A-Z=9.
\end{align*}
}
\end{vidu}

